\documentclass[a4paper, twoside]{report}

\usepackage[english]{babel}
\usepackage[utf8x]{inputenc}
\usepackage[T1]{fontenc}
\usepackage{listings}
\usepackage{hyperref}
\hypersetup{colorlinks=false}
\usepackage{lscape}
\usepackage{subfigure}
\usepackage{amsmath}
\usepackage{graphicx}
\usepackage[colorinlistoftodos]{todonotes}

%% Sets page size and margins
\usepackage[a4paper,top=3cm,bottom=2cm,left=3cm,right=3cm,marginparwidth=1.75cm]{geometry}

\title{Project Title}
\author{John Smith}
% Update supervisor and other title stuff in title/title.tex

\begin{document}
\begin{titlepage}

\newcommand{\HRule}{\rule{\linewidth}{0.5mm}} % Defines a new command for the horizontal lines, change thickness here

%----------------------------------------------------------------------------------------
%	LOGO SECTION
%----------------------------------------------------------------------------------------

\includegraphics[width=8cm]{title/logo.png}\\[1cm] % Include a department/university logo - this will require the graphicx package
 
%----------------------------------------------------------------------------------------

\center % Center everything on the page

%----------------------------------------------------------------------------------------
%	HEADING SECTIONS
%----------------------------------------------------------------------------------------
\quad\\[1.5cm]
%\textsc{\LARGE MSc Thesis}\\[1.5cm] % Name of your university/college
\textsc{\Large Imperial College London}\\[0.5cm] % Major heading such as course name
\textsc{\large Department of Electrical and Electronic Engineering}\\[0.5cm] % Minor heading such as course title

%----------------------------------------------------------------------------------------
%	TITLE SECTION
%----------------------------------------------------------------------------------------
\makeatletter
\HRule \\[0.4cm]
{ \huge \bfseries \@title}\\[0.4cm] % Title of your document
\HRule \\[1.5cm]
 
%----------------------------------------------------------------------------------------
%	AUTHOR SECTION
%----------------------------------------------------------------------------------------

\begin{minipage}{0.4\textwidth}
\begin{flushleft} \large
\emph{Author:}\\
\@author % Your name
\end{flushleft}
\end{minipage}
~
\begin{minipage}{0.4\textwidth}
\begin{flushright} \large
\emph{Supervisor:} \\
Prof. Lily Smith 
% Uncomment the following lines if there's a co-supervisor
%\\[1.2em] % Supervisor's Name
%\emph{Co-Supervisor:} \\
%Dr. Adam Smith % second marker's name
\end{flushright}
\end{minipage}\\[3cm]
\makeatother


%----------------------------------------------------------------------------------------
%	DATE SECTION
%----------------------------------------------------------------------------------------

{\large A thesis submitted for the degree of}\\[0.5cm]
{\large \emph{MSc Future Power Networks}}\\[0.5cm]
{\large \today}\\[2cm] % Date, change the \today to a set date if you want to be precise

\vfill % Fill the rest of the page with whitespace

\end{titlepage}

\begin{abstract}
Your abstract goes here. The abstract is a very brief summary of the dissertation's contents. It should be about half a page long. Somebody unfamiliar with your project should have a good idea of what it's about having read the abstract alone and will know whether it will be of interest to them.
\end{abstract}

\renewcommand{\abstractname}{Acknowledgements}
\begin{abstract}
It is usual to thank those individuals who have provided particularly useful assistance, technical or otherwise, during your project.
\end{abstract}

\tableofcontents
\listoffigures
\listoftables

\chapter{Introduction}

%\section{Background}\label{background}
%  Energy-constrained wireless devices are conventionally powered by batteries. However, the development of large-scale networks as Internet-of-Things (IoT) is strictly restricted by its limited working time and frequent recharging or replacement. Although Wireless Power Transfer (WPT) via inductive coupling has enjoyed some success in real-world applications, it is impractical for most devices on the move since the operation range is very short. As a promising alternative, the Radio-Frequency (RF) wave is typically with lower power level (\si{\uW} to \si{W}) but larger coverage (up to hundreds of meters) \cite{Ng2019}. Interestingly, it indeed carries both information and energy simultaneously, with the potential to power billions of mobile nodes wirelessly while keeping them connected. The recent revolution in harvester model and the significant power drop of electronics bring more possibility to the research on Wireless Information and Power Transfer (WIPT) via RF signals.


\section{Literature Review}\label{literature-review}
  Battery has become a main power source for most mobile devices. However, limited operation time and high cost in recharging or replacement have become bottlenecks for smart networks as Internet-of-Things (IoT). As a promising solution, Energy Harvesting (EH) from the ambient environment can potentially provide perpetual power to the devices. Compared with other renewable resources as solar, wind and water, Radio-Frequency (RF) waves typically contain less energy and are more suitable for low-power applications as Wireless Sensor Network (WSN). Recently, the significant reduction in power requirements of chips and processors brings more attention to Wireless Power Transfer (WPT) in both academia and industry \cite{R.Varshney2008, Grover2010, Zhang2013, Hui2014, Krikidis2014, Valenta2014, Boshkovska2015, Ding2015, Costanzo2016, Clerckx2018a}.

On the other hand, RF radiation has been a medium for Wireless Information Transfer (WIT) for more than a century. Naturally, a unified design of Wireless Information and Power transmission (WIPT) is expected to be a prominent solution to power billions of mobile devices while keeping them connected. \cite{R.Varshney2008} first defined a nonlinear concave capacity-energy function and investigated the tradeoff for typical binary channels and a flat additive white Gaussian noise (AWGN) channel with amplitude-constrained inputs. It was extended to frequency-selective channel in \cite{Grover2010}. However, both works were based on the impractical assumption that information decoding (ID) and EH can be performed individually on the same received signal. In \cite{Zhang2013}, the authors proposed two practical co-located receiver designs named \textit{time switching} (TS) that switches between ID and EH and \textit{power splitting} (PS) that splits the received power into two separate streams. It was then demonstrated in \cite{Zhou2013a} that TS can guarantee the same rate as conventional Time-Division Multiple Access (TDMA) while providing considerable power. In comparison, PS may lead to higher rate when the demand on power is sufficiently high. A further research \cite{Liu2013} enabled dynamic power splitting that adjusts the power split ratio based on the channel state information (CSI), and proposed a suboptimal low-complexity \textit{antenna switching} scheme. Nevertheless, the literatures above are mostly based on an oversimplified linear harvester model. To accurately characterize the behavior of the rectenna, \cite{Clerckx2016} derived a tractable nonlinear model and performed an adaptive multisine waveform design accordingly. Realistic simulations showed significant gains in harvested power and stressed the importance of modelling rectifier nonlinearity in wireless system design. The work was extended to multi-input single-output (MISO) WIPT in \cite{Clerckx2018} where a superposition of multicarrier modulated and unmodulated waveform was optimized as a function of CSI under transmit power budget. It suggested the rectifier nonlinearity can lead to a larger rate-energy (R-E) region and favours a different waveform, modulation and input distribution. In another perspective, a learning approach \cite{Varasteh2018} modelled the transmitter and receiver as deep neural networks (NN) and jointly optimized signal encoding with network parameters. Constellations showed that the offset of the power symbol is positively correlated to the power demand, while the information symbols are symmetrically located around the origin. The pattern confirmed a unmodulated waveform is beneficial to increase the harvested power in \cite{Clerckx2018}.
 
Energy-constrained wireless devices are conventionally powered by batteries. However, the development of large-scale networks as Internet-of-Things (IoT) is strictly restricted by its limited working time and frequent recharging or replacement. Although Wireless Power Transfer (WPT) via inductive coupling has enjoyed some success in real-world applications, it is impractical for most devices on the move since the operation range is very short. As a promising alternative, the Radio-Frequency (RF) wave is typically with lower power level (\si{\uW} to \si{W}) but larger coverage (up to hundreds of meters) \cite{Ng2019}. Interestingly, it indeed carries both information and energy simultaneously, with the potential to power billions of mobile nodes wirelessly while keeping them connected. The recent revolution in harvester model and the significant power drop of electronics bring more possibility to the research on Wireless Information and Power Transfer (WIPT) via RF signals.

\chapter{Main sections of the project}
\chapter{Evaluation}
\chapter{Conclusion}
\appendix
\chapter{Appendix}
The source code and relevant materials can be retrieved from \url{https://github.com/SnowzTail/}.

\bibliographystyle{unsrt}
\bibliography{bibs/sample}

\end{document}